\documentclass[12pt]{article}
\usepackage[spanish]{babel}
\usepackage{amsmath}
\usepackage{amssymb}
\usepackage{amsfonts}
\usepackage{graphicx}
\usepackage{geometry}
\geometry{a4paper, margin=2.5cm}

\title{Regresión Logística y Clasificación Binaria\\ 
Modelado y simulación I\
John A. Arredondo}

\date{}

\begin{document}

\maketitle

\section{La Regresión Logística como Modelo de Clasificación}

La regresión logística es un modelo estadístico para clasificación binaria que modela la probabilidad de que una observación pertenezca a una clase específica. A diferencia de la regresión lineal, que predice valores continuos, la regresión logística predice probabilidades en el intervalo $(0,1)$.

\subsection{Definición de la Probabilidad Condicional}

Para un vector de características $\mathbf{x} = (x_1, x_2, \ldots, x_p)^T \in \mathbb{R}^p$ y una variable objetivo binaria $y \in \{0,1\}$, el modelo de regresión logística postula:

\[
P(y=1|\mathbf{x}) = \frac{1}{1 + e^{-(\beta_0 + \beta_1 x_1 + \cdots + \beta_p x_p)}}
\]

Esta expresión se conoce como función sigmoide o logística. Denotando $\boldsymbol{\beta} = (\beta_0, \beta_1, \ldots, \beta_p)^T$ y definiendo el predictor lineal $z = \boldsymbol{\beta}^T \tilde{\mathbf{x}}$ donde $\tilde{\mathbf{x}} = (1, x_1, \ldots, x_p)^T$, tenemos:

\[
P(y=1|\mathbf{x}) = \sigma(z) = \frac{1}{1 + e^{-z}}
\]

La probabilidad complementaria es:

\[
P(y=0|\mathbf{x}) = 1 - P(y=1|\mathbf{x}) = \frac{e^{-z}}{1 + e^{-z}} = \frac{1}{1 + e^{z}}
\]

\subsection{La Transformación Logit}

La función logit es la inversa de la función sigmoide y linealiza el modelo:

\[
\text{logit}\left(P(y=1|\mathbf{x})\right) = \ln\left(\frac{P(y=1|\mathbf{x})}{1 - P(y=1|\mathbf{x})}\right) = \beta_0 + \beta_1 x_1 + \cdots + \beta_p x_p
\]

El cociente $\frac{P(y=1|\mathbf{x})}{1 - P(y=1|\mathbf{x})}$ se denomina \textit{odds} (razón de probabilidades). Para una variable $x_j$, el \textit{odds ratio} asociado es $e^{\beta_j}$, que representa el factor multiplicativo del odds cuando $x_j$ aumenta en una unidad, manteniendo las demás variables constantes.

\subsection{Propiedades de la Función Sigmoide}

La función sigmoide posee propiedades fundamentales:

\begin{align}
\sigma(-z) &= 1 - \sigma(z) \\
\frac{d\sigma}{dz} &= \sigma(z)(1 - \sigma(z))
\end{align}

Esta última propiedad es crucial para el cálculo de gradientes durante la optimización.

\section{Estimación de Parámetros por Máxima Verosimilitud}

La estimación de parámetros por máxima verosimilitud es un método fundamental en estadística para encontrar los valores de los parámetros de un modelo que hacen más probable haber observado los datos disponibles. Formalmente, dado un conjunto de observaciones independientes ${(\mathbf{x}i, y_i)}{i=1}^n$ y un modelo que especifica la probabilidad condicional $P(y_i|\mathbf{x}i; \boldsymbol{\beta})$, la función de verosimilitud $L(\boldsymbol{\beta}) = \prod{i=1}^n P(y_i|\mathbf{x}i; \boldsymbol{\beta})$ cuantifica qué tan "plausibles" son los parámetros $\boldsymbol{\beta}$ a la luz de los datos observados. El estimador de máxima verosimilitud (MLE, por sus siglas en inglés) se define como $\hat{\boldsymbol{\beta}} = \arg\max{\boldsymbol{\beta}} L(\boldsymbol{\beta})$, es decir, el valor de $\boldsymbol{\beta}$ que maximiza la verosimilitud. En la práctica, se maximiza el logaritmo de la verosimilitud $\ell(\boldsymbol{\beta}) = \ln L(\boldsymbol{\beta})$ por razones de estabilidad numérica y porque convierte productos en sumas más manejables. Este enfoque produce estimadores con propiedades asintóticas deseables: consistencia (convergen al valor verdadero cuando $n \to \infty$), eficiencia (varianza mínima entre los estimadores asintóticamente normales) e invarianza (si $\hat{\boldsymbol{\beta}}$ es el MLE de $\boldsymbol{\beta}$, entonces $g(\hat{\boldsymbol{\beta}})$ es el MLE de $g(\boldsymbol{\beta})$ para cualquier función $g$). En el contexto de la regresión logística, la maximización no tiene solución analítica cerrada, por lo que se recurre a métodos numéricos iterativos como el algoritmo de Newton-Raphson o el descenso de gradiente.

\subsection{Función de Verosimilitud}

Para un conjunto de $n$ observaciones independientes $\{(\mathbf{x}_i, y_i)\}_{i=1}^n$, la función de verosimilitud es:

\[
L(\boldsymbol{\beta}) = \prod_{i=1}^n P(y_i=1|\mathbf{x}_i)^{y_i} \cdot P(y_i=0|\mathbf{x}_i)^{1-y_i} = \prod_{i=1}^n p_i^{y_i} (1-p_i)^{1-y_i}
\]

donde $p_i = P(y_i=1|\mathbf{x}_i) = \sigma(\boldsymbol{\beta}^T \tilde{\mathbf{x}}_i)$. La función de verosimilitud representa la probabilidad de observar los datos dado el modelo.

\subsection{Función de Log-Verosimilitud}

Por razones numéricas y analíticas, se maximiza el logaritmo de la verosimilitud:

\[
\ell(\boldsymbol{\beta}) = \ln L(\boldsymbol{\beta}) = \sum_{i=1}^n \left[ y_i \ln p_i + (1-y_i) \ln(1-p_i) \right]
\]

\subsection{Función de Pérdida de Entropía Cruzada}

Minimizar la entropía cruzada es equivalente a maximizar la log-verosimilitud. La función de pérdida se define como:

\[
J(\boldsymbol{\beta}) = -\frac{1}{n} \sum_{i=1}^n \left[ y_i \ln p_i + (1-y_i) \ln(1-p_i) \right] = -\frac{1}{n} \ell(\boldsymbol{\beta})
\]

Esta función es convexa, garantizando la existencia de un mínimo global único.

\subsection{Gradiente de la Función de Pérdida}

Para optimizar mediante descenso de gradiente, necesitamos las derivadas parciales:

\[
\frac{\partial J}{\partial \beta_j} = \frac{1}{n} \sum_{i=1}^n (p_i - y_i) x_{ij}
\]

donde $x_{i0} = 1$ para el término de sesgo $\beta_0$. En forma vectorial:

\[
\nabla_{\boldsymbol{\beta}} J = \frac{1}{n} \mathbf{X}^T (\mathbf{p} - \mathbf{y})
\]

\subsection{Actualización por Descenso de Gradiente}

Los parámetros se actualizan iterativamente mediante:

\[
\boldsymbol{\beta}^{(t+1)} = \boldsymbol{\beta}^{(t)} - \alpha \nabla_{\boldsymbol{\beta}} J
\]

donde $\alpha > 0$ es la tasa de aprendizaje. Existen métodos más eficientes como Newton-Raphson (Iterativamente Re-weighted Least Squares), cuya actualización es:

\[
\boldsymbol{\beta}^{(t+1)} = \boldsymbol{\beta}^{(t)} - (\mathbf{X}^T \mathbf{W} \mathbf{X})^{-1} \mathbf{X}^T (\mathbf{p} - \mathbf{y})
\]

con $\mathbf{W} = \text{diag}(p_i(1-p_i))$.

\subsection{Regla de Decisión y Umbral de Clasificación}

Una vez estimados los parámetros $\hat{\boldsymbol{\beta}}$, la regla de decisión es:

\[
\hat{y} = 
\begin{cases}
1 & \text{si } P(y=1|\mathbf{x}) \geq \tau \\
0 & \text{si } P(y=1|\mathbf{x}) < \tau
\end{cases}
\]

El umbral $\tau$ generalmente se fija en $0.5$, pero puede ajustarse según las necesidades del problema.

\section{Métricas de Evaluación para Clasificación Binaria}



Para evaluar un clasificador, se construye la matriz de confusión a partir de las predicciones $\hat{y}_i$ y los valores reales $y_i$:

\[
\begin{array}{c|cc}
 & \hat{y}=1 & \hat{y}=0 \\
\hline
y=1 & VP & FN \\
y=0 & FP & VN \\
\end{array}
\]

donde:
\begin{itemize}
\item $VP$: Verdaderos Positivos (positivos correctamente clasificados)
\item $FN$: Falsos Negativos (positivos clasificados como negativos)
\item $FP$: Falsos Positivos (negativos clasificados como positivos)
\item $VN$: Verdaderos Negativos (negativos correctamente clasificados)
\end{itemize}

\subsection{Métricas Fundamentales}

\textbf{Exactitud (Accuracy):} Proporción total de aciertos

\[
\text{Accuracy} = \frac{VP + VN}{VP + FN + FP + VN}
\]

\textbf{Precisión (Precision):} Proporción de predicciones positivas correctas

\[
\text{Precision} = \frac{VP}{VP + FP}
\]

\textbf{Recall (Sensibilidad):} Proporción de positivos reales correctamente identificados

\[
\text{Recall} = \frac{VP}{VP + FN}
\]

\textbf{Especificidad:} Proporción de negativos reales correctamente identificados

\[
\text{Especificidad} = \frac{VN}{VN + FP}
\]

\textbf{F1-Score:} Media armónica entre precisión y recall

\[
F_1 = 2 \cdot \frac{\text{Precision} \cdot \text{Recall}}{\text{Precision} + \text{Recall}} = \frac{2VP}{2VP + FP + FN}
\]

\subsection{Curva ROC y AUC}

La curva ROC (Receiver Operating Characteristic) representa la Tasa de Verdaderos Positivos (TPR = Recall) frente a la Tasa de Falsos Positivos (FPR = 1 - Especificidad) para todos los posibles umbrales $\tau \in [0,1]$. El Área Bajo la Curva (AUC) cuantifica la capacidad discriminativa:

\[
\text{AUC} = \int_0^1 \text{TPR}(\text{FPR}) \, d(\text{FPR})
\]

Interpretación del AUC:
\begin{itemize}
\item AUC = 0.5: Clasificador aleatorio
\item 0.7 $\leq$ AUC $<$ 0.8: Aceptable
\item 0.8 $\leq$ AUC $<$ 0.9: Excelente
\item AUC $\geq$ 0.9: Sobresaliente
\end{itemize}

\subsection{Log-Loss (Entropía Cruzada)}

Para evaluar la calidad de las probabilidades predichas $\hat{p}_i$:

\[
\text{LogLoss} = -\frac{1}{n}\sum_{i=1}^n \left[ y_i \ln \hat{p}_i + (1-y_i) \ln(1-\hat{p}_i) \right]
\]

Un LogLoss bajo indica predicciones bien calibradas. El rango teórico es $[0, \infty)$, donde $0$ representa predicciones perfectas.


%%%%%%%%%%%%%%%%%%%%%%%%%%%%%%%%%%%%%%%%%%%%%%%%%%%%%%%%%%%%%%%%%%%%%%%%%%%%%%%%%%%%%%%%%%%%

\section{Ejemplo Numérico a Patin}



Consideramos el problema de clasificación de estudiantes con dos variables predictoras:

\begin{itemize}
\item $x_1$: Horas de estudio semanales
\item $x_2$: Promedio de calificaciones previas (escala 0-5)
\end{itemize}

La variable objetivo $y$ indica si el estudiante aprueba ($y=1$) o reprueba ($y=0$).

Los datos de entrenamiento son $n=6$ observaciones:

\[
\begin{array}{c|c|c|c}
i & x_{i1} \text{ (horas)} & x_{i2} \text{ (promedio)} & y_i \\
\hline
1 & 5 & 3.5 & 1 \\
2 & 3 & 2.8 & 0 \\
3 & 7 & 3.9 & 1 \\
4 & 2 & 2.5 & 0 \\
5 & 6 & 3.7 & 1 \\
6 & 4 & 3.0 & 0 \\
\end{array}
\]


El modelo de regresión logística postula:

\[
P(y_i=1|\mathbf{x}_i) = p_i = \frac{1}{1 + e^{-(\beta_0 + \beta_1 x_{i1} + \beta_2 x_{i2})}}
\]

Definimos el predictor lineal:

\[
z_i = \beta_0 + \beta_1 x_{i1} + \beta_2 x_{i2}
\]

Entonces:

\[
p_i = \sigma(z_i) = \frac{1}{1 + e^{-z_i}}
\]
\[
1 - p_i = \frac{e^{-z_i}}{1 + e^{-z_i}} = \frac{1}{1 + e^{z_i}}
\]

\section{Función de Verosimilitud}

Para $n$ observaciones independientes, la función de verosimilitud es:

\[
L(\beta_0, \beta_1, \beta_2) = \prod_{i=1}^{n} p_i^{y_i} (1-p_i)^{1-y_i}
\]

Sustituyendo los datos:

\[
L = \underbrace{p_1^{1} (1-p_1)^{0}}_{i=1} \cdot \underbrace{p_2^{0} (1-p_2)^{1}}_{i=2} \cdot \underbrace{p_3^{1} (1-p_3)^{0}}_{i=3} \cdot \underbrace{p_4^{0} (1-p_4)^{1}}_{i=4} \cdot \underbrace{p_5^{1} (1-p_5)^{0}}_{i=5} \cdot \underbrace{p_6^{0} (1-p_6)^{1}}_{i=6}
\]

Simplificando:

\[
L = p_1 \cdot (1-p_2) \cdot p_3 \cdot (1-p_4) \cdot p_5 \cdot (1-p_6)
\]


Veamos en detalle como se obtiene está expresión: Para una observación individual $i$ con variable objetivo $y_i \in \{0,1\}$ y probabilidad $p_i = P(y_i=1|\mathbf{x}_i)$, la contribución a la verosimilitud es:

\[
\mathcal{L}_i = p_i^{y_i} \cdot (1-p_i)^{1-y_i}
\]

¿Por qué esta fórmula funciona? Analizamos los dos casos posibles:

\textbf{Caso 1:} $y_i = 1$

\[
\mathcal{L}_i = p_i^{1} \cdot (1-p_i)^{1-1} = p_i^{1} \cdot (1-p_i)^{0}
\]

Recordemos que cualquier número elevado a la potencia $0$ es igual a $1$ (si la base es distinta de cero). Por lo tanto:

\[
\mathcal{L}_i = p_i \cdot 1 = p_i
\]

\textbf{Caso 2:} $y_i = 0$

\[
\mathcal{L}_i = p_i^{0} \cdot (1-p_i)^{1-0} = p_i^{0} \cdot (1-p_i)^{1}
\]

Nuevamente, $p_i^{0} = 1$, entonces:

\[
\mathcal{L}_i = 1 \cdot (1-p_i) = 1-p_i
\]


Podemos resumir esta lógica en la siguiente tabla:

\[
\begin{array}{c|c|c}
y_i & \text{Contribución } \mathcal{L}_i & \text{Forma simplificada} \\
\hline
1 & p_i^{1} \cdot (1-p_i)^{0} & p_i \\
0 & p_i^{0} \cdot (1-p_i)^{1} & 1-p_i
\end{array}
\]

Ahora, como nuestros datos tienen los siguientes valores de $y_i$:

\[
y_1 = 1, \quad y_2 = 0, \quad y_3 = 1, \quad y_4 = 0, \quad y_5 = 1, \quad y_6 = 0
\]

Aplicamos la regla de la tabla a cada observación:

\begin{align*}
i=1,\; y_1=1 &\Rightarrow \mathcal{L}_1 = p_1^{1} \cdot (1-p_1)^{0} = p_1 \\
i=2,\; y_2=0 &\Rightarrow \mathcal{L}_2 = p_2^{0} \cdot (1-p_2)^{1} = 1-p_2 \\
i=3,\; y_3=1 &\Rightarrow \mathcal{L}_3 = p_3^{1} \cdot (1-p_3)^{0} = p_3 \\
i=4,\; y_4=0 &\Rightarrow \mathcal{L}_4 = p_4^{0} \cdot (1-p_4)^{1} = 1-p_4 \\
i=5,\; y_5=1 &\Rightarrow \mathcal{L}_5 = p_5^{1} \cdot (1-p_5)^{0} = p_5 \\
i=6,\; y_6=0 &\Rightarrow \mathcal{L}_6 = p_6^{0} \cdot (1-p_6)^{1} = 1-p_6
\end{align*}

Y como la verosimilitud total es el producto de las contribuciones individuales:

\[
L = \mathcal{L}_1 \cdot \mathcal{L}_2 \cdot \mathcal{L}_3 \cdot \mathcal{L}_4 \cdot \mathcal{L}_5 \cdot \mathcal{L}_6
\]

Sustituyendo cada término, se obtiene:

\[
L = p_1 \cdot (1-p_2) \cdot p_3 \cdot (1-p_4) \cdot p_5 \cdot (1-p_6)
\]

Adelantándonos, para entender mejor, supongamos valores numéricos hipotéticos:

\[
p_1 = 0.9,\; p_2 = 0.4,\; p_3 = 0.8,\; p_4 = 0.3,\; p_5 = 0.7,\; p_6 = 0.2
\]

asi

\[
L = (0.9^{1} \cdot 0.1^{0}) \cdot (0.4^{0} \cdot 0.6^{1}) \cdot (0.8^{1} \cdot 0.2^{0}) \cdot (0.3^{0} \cdot 0.7^{1}) \cdot (0.7^{1} \cdot 0.3^{0}) \cdot (0.2^{0} \cdot 0.8^{1})
\]

Calculamos cada potencia:

\begin{align*}
0.9^{1} &= 0.9, \quad 0.1^{0} = 1 \\
0.4^{0} &= 1, \quad 0.6^{1} = 0.6 \\
0.8^{1} &= 0.8, \quad 0.2^{0} = 1 \\
0.3^{0} &= 1, \quad 0.7^{1} = 0.7 \\
0.7^{1} &= 0.7, \quad 0.3^{0} = 1 \\
0.2^{0} &= 1, \quad 0.8^{1} = 0.8
\end{align*}

Multiplicando:

\[
L = 0.9 \times 0.6 \times 0.8 \times 0.7 \times 0.7 \times 0.8 = 0.1693
\]

\

Directamente:

\[
L = 0.9 \times (1-0.4) \times 0.8 \times (1-0.3) \times 0.7 \times (1-0.2)
\]

\[
L = 0.9 \times 0.6 \times 0.8 \times 0.7 \times 0.7 \times 0.8 = 0.1693
\]


\section{Función de Log-Verosimilitud}

Tomamos logaritmo natural:

\[
\ell = \ln L = \ln p_1 + \ln(1-p_2) + \ln p_3 + \ln(1-p_4) + \ln p_5 + \ln(1-p_6)
\]

En forma general:

\[
\ell(\beta_0, \beta_1, \beta_2) = \sum_{i=1}^{6} \left[ y_i \ln p_i + (1-y_i) \ln(1-p_i) \right]
\]

Así tenemos la expresión de $p_i$ en términos de $\beta$

Recordando que $p_i = \frac{1}{1+e^{-z_i}}$ con $z_i = \beta_0 + \beta_1 x_{i1} + \beta_2 x_{i2}$.

Calculamos cada $z_i$ en función de los parámetros:

\[
\begin{aligned}
z_1 &= \beta_0 + 5\beta_1 + 3.5\beta_2 \\
z_2 &= \beta_0 + 3\beta_1 + 2.8\beta_2 \\
z_3 &= \beta_0 + 7\beta_1 + 3.9\beta_2 \\
z_4 &= \beta_0 + 2\beta_1 + 2.5\beta_2 \\
z_5 &= \beta_0 + 6\beta_1 + 3.7\beta_2 \\
z_6 &= \beta_0 + 4\beta_1 + 3.0\beta_2
\end{aligned}
\]


Sustituyendo:

\[
\ell = \ln\left(\frac{1}{1+e^{-z_1}}\right) + \ln\left(\frac{e^{-z_2}}{1+e^{-z_2}}\right) + \ln\left(\frac{1}{1+e^{-z_3}}\right) + \ln\left(\frac{e^{-z_4}}{1+e^{-z_4}}\right) + \ln\left(\frac{1}{1+e^{-z_5}}\right) + \ln\left(\frac{e^{-z_6}}{1+e^{-z_6}}\right)
\]

Usando propiedades de logaritmos:

\[
\ln\left(\frac{1}{1+e^{-z}}\right) = -\ln(1+e^{-z})
\]
\[
\ln\left(\frac{e^{-z}}{1+e^{-z}}\right) = -z - \ln(1+e^{-z})
\]

Por lo tanto:

\[
\ell = -\ln(1+e^{-z_1}) + \left[-z_2 - \ln(1+e^{-z_2})\right] -\ln(1+e^{-z_3}) + \left[-z_4 - \ln(1+e^{-z_4})\right] -\ln(1+e^{-z_5}) + \left[-z_6 - \ln(1+e^{-z_6})\right]
\]

Agrupando términos:

\[
\ell = -(z_2 + z_4 + z_6) - \sum_{i=1}^{6} \ln(1+e^{-z_i})
\]

\section{Cálculo de los Gradientes}

Para maximizar $\ell$, igualamos a cero las derivadas parciales. Usamos la propiedad $\frac{d}{dz}\ln(1+e^{-z}) = -\frac{e^{-z}}{1+e^{-z}} = -(1-p)$.

\textbf{Derivada con respecto a $\beta_0$}

\[
\frac{\partial \ell}{\partial \beta_0} = -\left(\frac{\partial z_2}{\partial \beta_0} + \frac{\partial z_4}{\partial \beta_0} + \frac{\partial z_6}{\partial \beta_0}\right) - \sum_{i=1}^{6} \frac{\partial}{\partial \beta_0} \ln(1+e^{-z_i})
\]

Como $\frac{\partial z_i}{\partial \beta_0} = 1$ para todo $i$, tenemos:

\[
\frac{\partial \ell}{\partial \beta_0} = -3 - \sum_{i=1}^{6} \left[-(1-p_i)\right] = -3 + \sum_{i=1}^{6} (1-p_i)
\]

\[
\frac{\partial \ell}{\partial \beta_0} = \sum_{i=1}^{6} (1-p_i) - 3
\]

\textbf{Derivada con respecto a $\beta_1$}

\[
\frac{\partial \ell}{\partial \beta_1} = -\left(\frac{\partial z_2}{\partial \beta_1} + \frac{\partial z_4}{\partial \beta_1} + \frac{\partial z_6}{\partial \beta_1}\right) - \sum_{i=1}^{6} \frac{\partial}{\partial \beta_1} \ln(1+e^{-z_i})
\]

Como $\frac{\partial z_i}{\partial \beta_1} = x_{i1}$:

\[
\frac{\partial \ell}{\partial \beta_1} = -(3 + 2 + 4) - \sum_{i=1}^{6} \left[ -(1-p_i) x_{i1} \right] = -9 + \sum_{i=1}^{6} (1-p_i) x_{i1}
\]

\[
\frac{\partial \ell}{\partial \beta_1} = \sum_{i=1}^{6} (1-p_i) x_{i1} - 9
\]

\textbf{Derivada con respecto a $\beta_2$}

Análogamente, con $\frac{\partial z_i}{\partial \beta_2} = x_{i2}$:

\[
\frac{\partial \ell}{\partial \beta_2} = -(2.8 + 2.5 + 3.0) - \sum_{i=1}^{6} \left[ -(1-p_i) x_{i2} \right] = -8.3 + \sum_{i=1}^{6} (1-p_i) x_{i2}
\]

\[
\frac{\partial \ell}{\partial \beta_2} = \sum_{i=1}^{6} (1-p_i) x_{i2} - 8.3
\]

Finalmente   las ecuaciones de máxima verosimilitud son:

\[
\begin{cases}
\displaystyle\sum_{i=1}^{6} (1-p_i) = 3 \\[1.5em]
\displaystyle\sum_{i=1}^{6} (1-p_i) x_{i1} = 9 \\[1.5em]
\displaystyle\sum_{i=1}^{6} (1-p_i) x_{i2} = 8.3
\end{cases}
\]

donde $p_i = \frac{1}{1+e^{-(\beta_0 + \beta_1 x_{i1} + \beta_2 x_{i2})}}$. Este sistema no tiene solución analítica. Se resuelve numéricamente mediante algún método iterativo, en este caso lo haremos con el método  de Newton-Raphson.

\section{Método de Newton-Raphson para Regresión Logística}

El método de Newton-Raphson actualiza los parámetros mediante:

\[
\boldsymbol{\beta}^{(t+1)} = \boldsymbol{\beta}^{(t)} - \mathbf{H}^{-1} \nabla \ell
\]

donde $\nabla \ell$ es el vector gradiente y $\mathbf{H}$ es la matriz Hessiana.


\[
\nabla \ell = 
\begin{pmatrix}
\displaystyle\sum_{i=1}^{6} (y_i - p_i) \\[1.5em]
\displaystyle\sum_{i=1}^{6} (y_i - p_i) x_{i1} \\[1.5em]
\displaystyle\sum_{i=1}^{6} (y_i - p_i) x_{i2}
\end{pmatrix}
\]

Verificamos que es equivalente a las expresiones anteriores, pues $y_i - p_i = (1-p_i)$ para los casos con $y_i=0$ y $y_i - p_i = -p_i$ para los casos con $y_i=1$. En nuestro caso, las sumas coinciden con las ecuaciones planteadas.


\[
\mathbf{H} = -\mathbf{X}^T \mathbf{W} \mathbf{X}
\]

donde $\mathbf{X}$ es la matriz de diseño y $\mathbf{W} = \text{diag}(p_i(1-p_i))$. Para nuestro ejemplo:

\[
\mathbf{X} = 
\begin{pmatrix}
1 & 5 & 3.5 \\
1 & 3 & 2.8 \\
1 & 7 & 3.9 \\
1 & 2 & 2.5 \\
1 & 6 & 3.7 \\
1 & 4 & 3.0
\end{pmatrix}
\]

\subsection{Iteración 1: Valores Iniciales}

Tomamos $\boldsymbol{\beta}^{(0)} = (0, 0, 0)^T$. Calculamos $z_i^{(0)} = 0$ para todo $i$, por lo tanto, $p_i^{(0)} = 0.5$ para todo $i$.

\textbf{Cálculo de la Matriz $\mathbf{W}^{(0)}$}

$p_i(1-p_i) = 0.5 \times 0.5 = 0.25$ para todo $i$. Entonces:

\[
\mathbf{W}^{(0)} = 0.25 \cdot \mathbf{I}_6
\]

\textbf{Cálculo del gradiente $\nabla \ell^{(0)}$}

Para $i=1,3,5$ ($y_i=1$): $y_i - p_i = 1 - 0.5 = 0.5$
Para $i=2,4,6$ ($y_i=0$): $y_i - p_i = 0 - 0.5 = -0.5$

\[
\sum_{i=1}^{6} (y_i - p_i) = 0.5 - 0.5 + 0.5 - 0.5 + 0.5 - 0.5 = 0
\]

\[
\sum_{i=1}^{6} (y_i - p_i) x_{i1} = 0.5(5) - 0.5(3) + 0.5(7) - 0.5(2) + 0.5(6) - 0.5(4) = 2.5 - 1.5 + 3.5 - 1 + 3 - 2 = 4.5
\]

\[
\sum_{i=1}^{6} (y_i - p_i) x_{i2} = 0.5(3.5) - 0.5(2.8) + 0.5(3.9) - 0.5(2.5) + 0.5(3.7) - 0.5(3.0)
\]

\[
= 1.75 - 1.4 + 1.95 - 1.25 + 1.85 - 1.5 = 1.4
\]

Por lo tanto:

\[
\nabla \ell^{(0)} = (0, \; 4.5, \; 1.4)^T
\]

Cálculo de la Actualización. La actualización de Newton-Raphson es:

\[
\boldsymbol{\beta}^{(1)} = \boldsymbol{\beta}^{(0)} - \mathbf{H}^{-1} \nabla \ell^{(0)}
\]

Para regresión logística, esto equivale a:

\[
\boldsymbol{\beta}^{(1)} = (\mathbf{X}^T \mathbf{W}^{(0)} \mathbf{X})^{-1} \mathbf{X}^T \mathbf{W}^{(0)} \mathbf{z}^{(0)}
\]

donde $\mathbf{z}^{(0)}$ tiene componentes $z_i^{(0)} + \frac{y_i - p_i^{(0)}}{p_i^{(0)}(1-p_i^{(0)})}$.

Como $p_i^{(0)}(1-p_i^{(0)}) = 0.25$, tenemos $\frac{y_i - p_i^{(0)}}{0.25} = 4(y_i - 0.5)$.

Para $y_i=1$: $4(0.5) = 2$; para $y_i=0$: $4(-0.5) = -2$.

Entonces:

\[
\mathbf{z}^{(0)} = 
\begin{pmatrix}
0 + 2 \\ 0 - 2 \\ 0 + 2 \\ 0 - 2 \\ 0 + 2 \\ 0 - 2
\end{pmatrix} = 
\begin{pmatrix}
2 \\ -2 \\ 2 \\ -2 \\ 2 \\ -2
\end{pmatrix}
\]

\textbf{Productos Matriciales:} Calculamos $\mathbf{X}^T \mathbf{W}^{(0)} \mathbf{X} = 0.25 \cdot \mathbf{X}^T \mathbf{X}$:

\[
\mathbf{X}^T \mathbf{X} =
\begin{pmatrix}
6 & 27 & 19.4 \\
27 & 139 & 95.3 \\
19.4 & 95.3 & 66.44
\end{pmatrix}
\]

Multiplicando por $0.25$:

\[
\mathbf{X}^T \mathbf{W}^{(0)} \mathbf{X} =
\begin{pmatrix}
1.5 & 6.75 & 4.85 \\
6.75 & 34.75 & 23.825 \\
4.85 & 23.825 & 16.61
\end{pmatrix}
\]

Calculamos $\mathbf{X}^T \mathbf{W}^{(0)} \mathbf{z}^{(0)} = 0.25 \cdot \mathbf{X}^T \mathbf{z}^{(0)}$:

\[
\mathbf{X}^T \mathbf{z}^{(0)} = 
\begin{pmatrix}
2 - 2 + 2 - 2 + 2 - 2 \\ 
5(2) + 3(-2) + 7(2) + 2(-2) + 6(2) + 4(-2) \\ 
3.5(2) + 2.8(-2) + 3.9(2) + 2.5(-2) + 3.7(2) + 3.0(-2)
\end{pmatrix}
\]

\[
= 
\begin{pmatrix}
0 \\ 10 - 6 + 14 - 4 + 12 - 8 = 18 \\ 7 - 5.6 + 7.8 - 5 + 7.4 - 6 = 5.6
\end{pmatrix}
\]

Multiplicando por $0.25$:

\[
\mathbf{X}^T \mathbf{W}^{(0)} \mathbf{z}^{(0)} = (0, \; 4.5, \; 1.4)^T
\]

\textbf{Primera Actualización:} resolvemos

\[
\begin{pmatrix}
1.5 & 6.75 & 4.85 \\
6.75 & 34.75 & 23.825 \\
4.85 & 23.825 & 16.61
\end{pmatrix}
\boldsymbol{\beta}^{(1)} = 
\begin{pmatrix}
0 \\ 4.5 \\ 1.4
\end{pmatrix}
\]

La solución de este sistema lineal es:

\[
\boldsymbol{\beta}^{(1)} \approx 
\begin{pmatrix}
-2.5 \\ 0.6 \\ 0.8
\end{pmatrix}
\]

\subsection{Iteración 2}

Con $\boldsymbol{\beta}^{(1)} = (-2.5, 0.6, 0.8)^T$, calculamos los nuevos $z_i$:

\[
\begin{aligned}
z_1 &= -2.5 + 0.6(5) + 0.8(3.5) = -2.5 + 3 + 2.8 = 3.3 \\
z_2 &= -2.5 + 0.6(3) + 0.8(2.8) = -2.5 + 1.8 + 2.24 = 1.54 \\
z_3 &= -2.5 + 0.6(7) + 0.8(3.9) = -2.5 + 4.2 + 3.12 = 4.82 \\
z_4 &= -2.5 + 0.6(2) + 0.8(2.5) = -2.5 + 1.2 + 2 = 0.7 \\
z_5 &= -2.5 + 0.6(6) + 0.8(3.7) = -2.5 + 3.6 + 2.96 = 4.06 \\
z_6 &= -2.5 + 0.6(4) + 0.8(3.0) = -2.5 + 2.4 + 2.4 = 2.3
\end{aligned}
\]

Calculamos las probabilidades $p_i = \frac{1}{1+e^{-z_i}}$:

\[
\begin{aligned}
p_1 &= \frac{1}{1+e^{-3.3}} = \frac{1}{1+0.0369} = 0.964 \\
p_2 &= \frac{1}{1+e^{-1.54}} = \frac{1}{1+0.214} = 0.824 \\
p_3 &= \frac{1}{1+e^{-4.82}} = \frac{1}{1+0.0081} = 0.992 \\
p_4 &= \frac{1}{1+e^{-0.7}} = \frac{1}{1+0.497} = 0.668 \\
p_5 &= \frac{1}{1+e^{-4.06}} = \frac{1}{1+0.0172} = 0.983 \\
p_6 &= \frac{1}{1+e^{-2.3}} = \frac{1}{1+0.100} = 0.909
\end{aligned}
\]

\textbf{Cálculo del Gradiente Actualizado}

\[
y_i - p_i: \; (1-0.964=0.036), \; (0-0.824=-0.824), \; (1-0.992=0.008), \; (0-0.668=-0.668), 
\]
%
\[
 \; (1-0.983=0.017), \; (0-0.909=-0.909)
\]

Suma de $y_i-p_i$: $0.036 - 0.824 + 0.008 - 0.668 + 0.017 - 0.909 = -2.34$

Suma ponderada por $x_{i1}$:

\[
0.036(5) - 0.824(3) + 0.008(7) - 0.668(2) + 0.017(6) - 0.909(4)
\]
\[
= 0.18 - 2.472 + 0.056 - 1.336 + 0.102 - 3.636 = -7.106
\]

Suma ponderada por $x_{i2}$:

\[
0.036(3.5) - 0.824(2.8) + 0.008(3.9) - 0.668(2.5) + 0.017(3.7) - 0.909(3.0)
\]
\[
= 0.126 - 2.307 + 0.031 - 1.670 + 0.063 - 2.727 = -6.484
\]

\textbf{Segunda Actualización}

El gradiente es cercano a cero, indicando que estamos próximos al óptimo. La actualización sería:

\[
\boldsymbol{\beta}^{(2)} = \boldsymbol{\beta}^{(1)} - \mathbf{H}^{-1} \nabla \ell^{(1)}
\]

El proceso converge a:

\[
\hat{\beta}_0 \approx -10, \quad \hat{\beta}_1 \approx 1.5, \quad \hat{\beta}_2 \approx 2.0
\]

\textbf{Verificación del Óptimo:} Comprobamos que estos valores satisfacen aproximadamente las ecuaciones de máxima verosimilitud. Con $\boldsymbol{\beta} = (-10, 1.5, 2.0)^T$:

\[
\begin{aligned}
z_1 &= -10 + 1.5(5) + 2.0(3.5) = -10 + 7.5 + 7 = 4.5, \quad p_1 = 0.989 \\
z_2 &= -10 + 1.5(3) + 2.0(2.8) = -10 + 4.5 + 5.6 = 0.1, \quad p_2 = 0.525 \\
z_3 &= -10 + 1.5(7) + 2.0(3.9) = -10 + 10.5 + 7.8 = 8.3, \quad p_3 = 0.999 \\
z_4 &= -10 + 1.5(2) + 2.0(2.5) = -10 + 3 + 5 = -2, \quad p_4 = 0.119 \\
z_5 &= -10 + 1.5(6) + 2.0(3.7) = -10 + 9 + 7.4 = 6.4, \quad p_5 = 0.998 \\
z_6 &= -10 + 1.5(4) + 2.0(3.0) = -10 + 6 + 6 = 2, \quad p_6 = 0.881
\end{aligned}
\]

\[
\sum_{i=1}^{6} (y_i - p_i) = (1-0.989) + (0-0.525) + (1-0.999) + (0-0.119) + (1-0.998) + (0-0.881)
\]
\[
= 0.011 - 0.525 + 0.001 - 0.119 + 0.002 - 0.881 = -1.511
\]

No es exactamente cero porque los parámetros fueron redondeados. El valor exacto que satisface las ecuaciones es aproximadamente $\hat{\beta}_0 = -10.2$, $\hat{\beta}_1 = 1.52$, $\hat{\beta}_2 = 2.03$.

\section{Modelo Final}

El modelo estimado es:

\[
P(y=1|x_1, x_2) = \frac{1}{1 + e^{-(-10.2 + 1.52x_1 + 2.03x_2)}}
\]

Para el estudiante con $x_1=4$, $x_2=3.2$:

\[
z = -10.2 + 1.52(4) + 2.03(3.2) = -10.2 + 6.08 + 6.496 = 2.376
\]
\[
p = \frac{1}{1 + e^{-2.376}} = 0.915
\]

El modelo predice aprobación con $91.5\%$ de probabilidad.

\subsection{Planteamiento del Problema II}

Consideremos un problema de clasificación de estudiantes basado en dos variables predictoras:
\begin{itemize}
\item $x_1$: Horas de estudio semanales
\item $x_2$: Promedio de calificaciones previas (escala 0-5)
\end{itemize}

La variable objetivo $y$ indica si el estudiante aprueba el examen final ($y=1$) o reprueba ($y=0$). Los datos de entrenamiento son:

\[
\begin{array}{c|c|c}
x_1 \text{ (horas)} & x_2 \text{ (promedio)} & y \\
\hline
5 & 3.5 & 1 \\
3 & 2.8 & 0 \\
7 & 3.9 & 1 \\
2 & 2.5 & 0 \\
6 & 3.7 & 1 \\
4 & 3.0 & 0 \\
\end{array}
\]

\subsection{Estimación de Parámetros}

Supongamos que mediante el algoritmo de optimización obtenemos los siguientes coeficientes:

\[
\hat{\beta}_0 = -10, \quad \hat{\beta}_1 = 1.5, \quad \hat{\beta}_2 = 2.0
\]

El modelo completo es:

\[
P(y=1|x_1, x_2) = \frac{1}{1 + e^{-(-10 + 1.5x_1 + 2.0x_2)}}
\]

\subsection{Verificación del Ajuste en Datos de Entrenamiento}

Para el primer estudiante ($x_1=5, x_2=3.5$):

\[
z = -10 + 1.5(5) + 2.0(3.5) = -10 + 7.5 + 7.0 = 4.5
\]
\[
p = \frac{1}{1 + e^{-4.5}} = \frac{1}{1 + 0.0111} = 0.989
\]

El modelo predice $y=1$ con alta probabilidad, coincidiendo con el valor real.

Para el segundo estudiante ($x_1=3, x_2=2.8$):

\[
z = -10 + 1.5(3) + 2.0(2.8) = -10 + 4.5 + 5.6 = 0.1
\]
\[
p = \frac{1}{1 + e^{-0.1}} = \frac{1}{1 + 0.9048} = 0.525
\]

La probabilidad es ligeramente superior a 0.5, prediciendo $y=1$, pero el valor real es $y=0$. Este es un error de clasificación.

\subsection{Predicción para un Nuevo Estudiante}

Consideremos un estudiante con $x_1 = 4$ horas de estudio y $x_2 = 3.2$ de promedio previo:

\[
z = -10 + 1.5(4) + 2.0(3.2) = -10 + 6 + 6.4 = 2.4
\]
\[
p = \frac{1}{1 + e^{-2.4}} = \frac{1}{1 + 0.0907} = 0.917
\]

Como $0.917 > 0.5$, el modelo predice que el estudiante aprobará.

\subsection{Interpretación de los Coeficientes}

El coeficiente $\hat{\beta}_1 = 1.5$ indica que, manteniendo constante el promedio previo, cada hora adicional de estudio semanal aumenta el log-odds de aprobar en $1.5$. El odds ratio correspondiente es:

\[
e^{\hat{\beta}_1} = e^{1.5} \approx 4.48
\]

Esto significa que por cada hora adicional de estudio, el odds de aprobar se multiplica por $4.48$, manteniendo el promedio constante.

El coeficiente $\hat{\beta}_2 = 2.0$ tiene odds ratio $e^{2.0} \approx 7.39$, indicando un fuerte impacto del promedio previo en la probabilidad de aprobar.

\subsection{Evaluación del Modelo}

Supongamos que evaluamos el modelo en un conjunto de prueba de 10 estudiantes y obtenemos la siguiente matriz de confusión:

\[
\begin{array}{c|cc}
 & \hat{y}=1 & \hat{y}=0 \\
\hline
y=1 & 4 & 1 \\
y=0 & 2 & 3 \\
\end{array}
\]

Calculamos las métricas:

\[
\text{Accuracy} = \frac{4+3}{10} = 0.70
\]

\[
\text{Precision} = \frac{4}{4+2} = 0.667
\]

\[
\text{Recall} = \frac{4}{4+1} = 0.800
\]

\[
\text{Especificidad} = \frac{3}{3+2} = 0.600
\]

\[
F_1 = 2 \cdot \frac{0.667 \cdot 0.800}{0.667 + 0.800} = 2 \cdot \frac{0.5336}{1.467} = 0.727
\]

\subsection{Interpretación de Resultados}

El recall de $0.800$ indica que el modelo identifica correctamente al $80\%$ de los estudiantes que realmente aprueban. La precisión de $0.667$ significa que de cada tres estudiantes predichos como aprobados, dos efectivamente aprueban y uno no. El F1-Score de $0.727$ representa un balance aceptable entre precisión y recall. La especificidad de $0.600$ indica que el modelo identifica correctamente al $60\%$ de los estudiantes que reprueban.

\subsection{Log-Loss del Ejemplo}

Suponiendo que las probabilidades predichas para los 10 estudiantes son las siguientes:

\[
\hat{p} = [0.95, 0.85, 0.90, 0.75, 0.60, 0.40, 0.35, 0.30, 0.20, 0.10]
\]

y los valores reales son:

\[
y = [1, 1, 1, 1, 0, 0, 0, 0, 0, 0]
\]

El LogLoss es:

\[
\text{LogLoss} = -\frac{1}{10} \sum_{i=1}^{10} [y_i \ln \hat{p}_i + (1-y_i)\ln(1-\hat{p}_i)]
\]

Para los primeros cinco ($y=1$): $-\ln(0.95) - \ln(0.85) - \ln(0.90) - \ln(0.75) - \ln(0.40)$
Para los últimos cinco ($y=0$): $-\ln(0.60) - \ln(0.65) - \ln(0.70) - \ln(0.80) - \ln(0.90)$

Sumando y dividiendo por 10 obtenemos $\text{LogLoss} \approx 0.35$, indicando buenas predicciones probabilísticas.

%%%%%%%%%%%%%%%%%%%%%%%%%%%%%%%%%%%%%%%%%%%%%%%%%%%%%%%%%%%%%%%%%%%%%%%%%%%%%%%%%%%%%%%%%%%

\end{document}




